\clearpage
\fancyhead[L]{\small BODE. ASTRONOMICAL YEARBOOK.}
\hypertarget{native-553-lower}{}\label{native-553-lower}
\begin{center}\small Göttingische gelehrte Anzeigen. Issue 22. Pages 209..220. 6 February 1813.\\[3pt]\rule{.72\linewidth}{.3pt}\end{center}

\emph{Astronomical Yearbook for the Year} 1814, together with a collection of the latest papers, observations, and reports in the astronomical sciences, calculated and edited by J. E. \textsc{Bode}, royal astronomer and member of the Academy of Sciences. 1811. From the author and on commission through J. E. \textsc{Hitzig}. Berlin. 275 pages in octavo and a copperplate.

Astronomical Yearbook for the Year 1815, etc. Same place, 1812. Also 275 pages in octavo and a copperplate.

Without repeating our remarks on the arrangement of the astronomical calendar from our reviews of earlier volumes of this yearbook, we confine ourselves here to reporting on the appended papers. The volume for 1814 contains the following 34 articles. 1) Mean astronomical refraction, according to \textsc{Laplace}. 2) Astronomical observations made by \textsc{Bugge} in 1809 and 1810 at the Royal Observatory in Copenhagen concern Ceres, Vesta, Uranus, Saturn, Jupiter, eclipses of Jupiter's satellites, and stellar occultations. 3) On a method of determining the time by measuring the distance of the Sun from a fixed and known point on the horizon, by Mr. J. F. \textsc{von Beek Calkoen}. This \SourceLeaf{554}method can yield usable results if due care is taken. Its use must extend beyond objects on the horizon, and due account must be taken of refraction and parallax. Both can be allowed for without making the procedure very lengthy. 4) Astronomical observations made at the Imperial Observatory in Vienna by Dr. \textsc{Triesnecker} and Professor \textsc{Bürg} include eclipses of satellites, stellar occultations, and observations of planets and the Sun. 5) The ideas on the calculation of perturbations according to \textsc{Kepler}, together with remarks, by Mr. J. W. \textsc{Pfaff}, deserve attention in part, although some of them strike us as immature or hasty. 6) Astronomical observations made in 1810 at the Royal Observatory in Prague by Professor \textsc{David} and Adjunct \textsc{Bittner}. We cannot approve of the way observed stellar altitudes are used here and in the preceding volumes of this yearbook for a supposed correction of the refraction tables. Anyone who wishes to work on this delicate point must first establish the altitude of the pole at the observing site more firmly, make far more observations overall than Mr. \textsc{David} has made, and must do more than borrow the stellar declinations from other astronomers. 7) From a letter from Professor \textsc{Littrow} in Kazan. It is gratifying to see from it how astronomy is flourishing on the border of Asia. 8) A more accurate determination of the period of variation in the brightness of the star $\eta$ Antinoi, by Professor \textsc{Wurm}. From his own observations, which span 24 years, Mr. \textsc{Wurm} finds a period of $7{,}1761$ days. One might have wished for a less arbitrary way of combining the observations, perhaps by the method of least squares, followed by a comparison with the individual observations. 9) Observations of the comet of 1807 and of the total solar eclipse of 16 June 1806 at Salem in the United States of North America, by Mr. \textsc{Bowditch}. From the latest volume of the \emph{Memoirs of the American Academy of Arts and Sciences}. The account of the circumstances accompanying the total solar eclipse is especially interesting. 10) A report on the Mannheim Observatory by State Councillor \textsc{Klüber}, curator of the observatory. What is said here is already known in greater detail from Mr. \textsc{Klüber}'s separate publication that has since appeared. 11) Astronomical observations made at the Royal Observatory in Berlin in 1810. We single out the determination of the altitude of the pole at the Berlin Observatory with the two-foot Troughton circle, for which Mr. \textsc{Bode} finds a mean of $52^\circ\,31'\,15''$. 12) On the comet of 1795, by Dr. \textsc{Olbers}. Observations of this comet are very scanty, and the elements obtained by different astronomers differ considerably. Dr. \textsc{Olbers} based a new determination of the elements partly on a previously unused observation by \textsc{Herschel} and partly on a newly reduced observation by Mr. \textsc{Bode}. 13) Results of an investigation into the orientation of the plane of Saturn's ring, the theory of its fourth satellite, the masses of the planet and the ring, and observed stellar occultations, by Professor \textsc{Bessel}. The detailed investigation of Saturn's fourth satellite has already appeared in the third issue of the Königsberg Archive. Its important results deserve the attention of every astronomer. 14) Calculation of the orbit of the comet of 1810, by the same author. 15) Astronomical observations made at Kremsmünster in 1809 and 1810 by Mr. \textsc{Derfflinger}. At the immersion of Aldebaran on 18 September 1810, the star seemed to Mr. \textsc{Derfflinger} to linger for a few seconds on the bright limb of the Moon with its light altered. \textsc{Bugge}, \textsc{Bode}, and \textsc{David} noticed the same thing, and some astronomers have regarded it as extraordinary. At our observatory, where the same occultation was also observed, nothing of the kind was noticed. We therefore leave open whether it was anything more than irradiation of the lunar limb seen by an eye that was perhaps already somewhat tired. 16) Observation and calculation of the occultation of Aldebaran by the Moon on 18 September 1810, \SourceLeaf{555}carried out in Dorpat by Professor \textsc{Knorre}. There is no mention here of the phenomenon just described. 17) On measuring heights with the barometer, by Professor \textsc{Benzenberg}. This attempt to present the subject for a general audience is commendable in itself, though it is perhaps a little out of place here. 18) Observations of Juno and Vesta in 1811 at the Imperial Observatory in Wilna, by Professor \textsc{Sniadecki}. 19) From a letter from Dr. \textsc{Pansner} in Petersburg. This article reports on a trigonometrical survey of the coast of the Gulf of Finland, whose completion is much to be desired. 20) On the accuracy of \textsc{Baumann}'s vertical circle, by Dr. \textsc{Pottgiesser} in Eberfeld. The paper shows considerable practical insight. Besides estimating the errors from the construction of the instrument itself, the author tested its accuracy with some actual observations, which speak very much in its favour. We would only have wished for a larger number and greater variety of observations, since recent experience has raised some doubts about circles with upright pillars. 21) Determinations of the longitude and latitude of several places in Austrian territory, together with observed stellar occultations, by Imperial Baroness von \textsc{Matt}. 22) Observations on the annual parallax of $\alpha$ of the Lyre, by Mr. \textsc{Calandrelli} in Rome. An annual parallax of $5''8$ is inferred here, though several objections could still be raised about its reliability. The sector used for the observations was left unreversed. Its collimation error was assumed to have remained unchanged throughout the year. 23) Design for a sundial that shows twelve noon in mean time. The figure-of-8 curve traced by the tip of the gnomon's shadow on a vertical sundial at mean noon in the various months of the year is drawn through 36 points. The use of mean time for public clocks certainly has much in its favour. However, the people responsible for setting clocks in the countryside or in smaller places are seldom properly instructed in the matter. The great advantage of \emph{agreement} among the public clocks is thus easily lost. This agreement matters particularly along postal routes and is more readily maintained when true time is used. 24) Astronomical determinations of position, by Chief Pastor \textsc{Fritsch} in Quedlinburg. On a trip to Silesia in the summer of 1810, Mr. \textsc{Fritsch} determined the altitude of the pole at Ballenstedt, Werningerode, Oschatz, Bischofswerda, and several places in Silesia. 25) On the comet of 1811 and its return in August, by Dr. \textsc{Olbers}. 26) Astronomical observations and remarks, by Professor \textsc{Gauss}, containing observations of Pallas and Vesta and corrections to some remarks by Dr. \textsc{Treisnecker} in the yearbook for 1813. The next three articles give the ephemerides for Pallas, Juno, and Vesta by Messrs. \textsc{Nicolai}, \textsc{Wachter}, and \textsc{Gerling}. 30) Observations and orbital elements of the comet of 1811, and a supplement to the Theoria motus corporum coelestium by Professor \textsc{Gauss}. We notice a misprint in the \emph{supplement}. On p. 257, line 5, $z$ should read $-z$. 31) Observations of the comet of 1811, the elements of its orbit, and stellar occultations, by Professor \textsc{Bessel}. 32) From a letter from Dr. \textsc{Koch} in Danzig. Remarks on some fixed stars and a determination of the declination of the magnetic needle in Danzig ($13^\circ\,48'$). 33) Observations of the comet of 1811 at the Royal Observatory in Berlin. It should be noted that the accompanying drawing of the comet's true orbit is not a projection. --- Among the brief astronomical notes that make up the last article, we single out only the report from a letter from Dr. B\dots\ (\textsc{Brinkley}?), who found with an eight-foot circle an annual parallax for $\alpha$ of the Lyre of $2''52$, which is less than half of \textsc{Calandrelli}'s result cited above.

The volume for 1815 contains 31 articles. It opens with remarks by Professor \textsc{Littrow} on Professor \SourceLeaf{556}\textsc{Gauss}'s Theoria motus corporum coelestium. They concern the problem of determining the orbit of a planet from three geocentric positions. For the admittedly very rare case in which three hypotheses do not suffice, Mr. \textsc{Littrow} makes a slight change in the course of the calculation for the fourth hypothesis. He uses the same device by which \textsc{Gauss} determines the new values of $P$ and $Q$ to calculate $r, r', f, f', f''$ at once. Although this procedure may sometimes offer some advantage, we would hesitate to recommend it without reservation. In general there is no reason to expect more accurate results from it than from the procedure used in the Theoria. Part of the saving is lost because it does not also give the orientation of the orbital plane, and applying the same device straightaway to that orientation as well can sometimes be rather problematic. We consider the next method, proposed by the same astronomer for determining a circular orbit from two geocentric positions, very well suited to its purpose, though we doubt that it is advantageous for a preliminary determination of the distances in \emph{cometary orbits}. We also find a straightforward derivation of aberration in longitude and latitude by the same author, as well as a pair of series expressions for refraction, one of which is insufficiently explained to allow a judgment. 2) Astronomical observations made in Pisa, Milan, and Padua (taken from volume XV of the Memorie di Matematica e di Fisica della Società Italiana). 3) Astronomical reports and remarks, physical observation of the great comet of 2811, geographical determinations, etc., by Professor \textsc{Huth} in Dorpat. 4) Observation of the great comet of 1811, calculation of its orbit, stellar occultations, etc., by Professor \textsc{Bessel}. 5) On the discovery of a new comet in November 1811, observations of it and of the great comet of 1811, and observation of Pallas, by Dr. \textsc{Olbers}. 6) Observations of the comet of 1807 and the elements of its orbit, by Mr. \textsc{Niccolo Cacciatore}, assistant at the Royal Observatory in Palermo. A very valuable supplement on the apparent orbit of that comet, as yet. unknown on the mainland. 7) Observed and calculated oppositions of Mars, Vesta, Jupiter, and Saturn, eclipses of Jupiter's satellites and stellar occultations in 1811, elements of the comet of 1810, observation of the great comet of 1811, and calculation of its orbit, by Dr. \textsc{Triesnecker}. 8) Observations of Uranus, Saturn, Mars, and Ceres, eclipses of Jupiter's satellites, stellar occultations, and observations of the great comet of 1811, made at the Imperial Observatory in Wilna by Professor \textsc{Sniadecki}. 9) Astronomical reports and observations, geographical determinations of position, etc., by Professor \textsc{Jabbo Oltmanns}. Stellar occultations observed in Amsterdam and Greenwich, observations in South America, and determinations of position in Spain and Portugal. 10) In remembrance of \textsc{Halley}'s comet, by Professor J. W. \textsc{Pfaff} in Nuremberg. Here the author gives himself over to the prospect of all that repeated observations of the return of comets will one day reveal. Admittedly, centuries must pass before even part of these hopes can be realised. 11) Astronomical observations made at the Royal Observatory in Berlin in 1811. 12) Observations concerning eclipses of Jupiter's satellites, lunar eclipses, stellar occultations, refraction, planetary oppositions, and the great comet of 1811, made at the Prague Observatory by Professor \textsc{David}. 13) Observations of Pallas and of the second comet of 1811, elements of the latter's orbit, stellar occultations, etc., by Professor \textsc{Gauss}. 14) On the Earth and the Moon coming to one and the same place. A calculation in advance, for the years 1812--1815, of the instances when the Earth arrives a few hours later at the place previously occupied by the Moon, and vice versa. As is well known, Lichtenberg once raised the question of whether striking changes in the weather might occur in the first of these cases. We admit that \SourceLeaf{557}we (like Mr. \textsc{Bode} himself) expect nothing of the kind, and would only add that the calculation itself rests on the inadmissible assumption that the solar system is absolutely at rest. 15) Observations of Ceres and of the oppositions of Uranus and Mars in 1811, made at the Kremsmünster Observatory by Mr. \textsc{Derfflinger}. 16) Astronomical remarks by Dr. \textsc{Lamberti} in Dorpat. 17) On the apparent orbit of the great comet of 1811, with a drawing. 18) On the motion of the double star 61 in Cygnus, by Professor \textsc{Bessel}. This remarkable subject has already been reported on in several other places, and first in our journal (25 April 1812, issue 67 [p. 349 of this volume]). 19) An improved determination of the altitude of the pole at Riga, by Professor \textsc{Sandt}. 20) Astronomical observations made in Paris and Greenwich in the years 1805--1809, communicated to Professor \textsc{Oltmanns}. 21) A report by Professor \textsc{Fischer} in Berlin on very fine parallel mirrors made by Mr. \textsc{Duve}, an instrument maker in Berlin. The importance of having perfectly flat and parallel surfaces on the glass mirrors in reflecting instruments makes Mr. \textsc{Duve}'s efforts, which according to Mr. \textsc{Fischer} have succeeded well, especially valuable. It is noteworthy that, as reported here, one of these parallel mirrors lost its former perfection after a small piece was accidentally broken off. 22) Astronomical observations and remarks, by Mr. Joseph \textsc{Bayr} at Kloster Hradisch near Ollmütz. Geographical determination of the position of Hradisch and altitude of the pole at Troppau. 23) On the influence of \textsc{Dalton}'s theory on the measurement of heights and on refraction, by Dr. \textsc{Benzenberg}. Judgment on \textsc{Dalton}'s ingenious but still too uncertain hypothesis belongs before the physicist’s tribunal and lies outside the astronomer’s province. We think chemistry should easily be able to decide whether Mr. \textsc{Benzenberg}'s reply to Mr. \textsc{Tralles}'s objections to that theory is admissible. For now it would be premature to take the hypothesis as an established truth on which to base barometric measurements of height. 24) Occasional thoughts on the surface of the Moon, by Lieutenant \textsc{von Boguslawski}. Although we see no water on the Moon's surface, in Mr. \textsc{von Boguslawski}'s opinion it could still be there, with the Sun's heat turning it on the dayside into invisible gaseous vapours. After sunset these would gradually turn into mist, dew, and ice as they lose their heat. The instantaneous disappearance of fixed stars when they are occulted by the dark lunar limb seems to us an important objection to this hypothesis. 25) Observation of Pallas and Juno, calculation of their oppositions, the elements of the latter's orbit, etc., by Professor \textsc{Gauss}. The two following articles give the ephemerides of Pallas and Juno for 1813 by Messrs. \textsc{Nicolai} and \textsc{Wachter}. 28) Further remarks on the variable double star 61 Cygnus, observation of the great comet of 1811, stellar occultations, and astronomical reports, by Professor \textsc{Bessel}. 29) Astronomical observations and remarks, by Dr. \textsc{Koch} in Danzig. 30) On the new comet of 1812. Under the heading “Miscellaneous Astronomical Reports”, the last article also gives Mr. \textsc{Schubert}'s observations of the great comet of 1811 in Petersburg, the prices of the astronomical instruments made by Mr. \textsc{Baumann} in Stuttgart, and various other astronomical notes. --- We close this review of a yearbook that remains as rich and interesting as ever by expressing our wish that its worthy editor will continue to receive strong support for his laborious undertaking and will for many years to come give us the pleasure of a new volume each year.
